\documentclass{article}
\title{Latent Policy Network}
\author{Ryan Marr}
\date{2026}

\begin{document}
\maketitle

\begin{abstract}
Latent Policy Network is a research framework for reinforcement learning policies
that generate their own weights on the fly. It targets rapidly changing 1v1
settings: the opponent has a hidden style that can switch mid-episode, and the
agent encodes its recent interaction history into a context that produces the
policy's weights, so it adapts to a new opponent without further training.
\end{abstract}

\begin{figure}
\includegraphics{imgs/latent-policy-network.gif}
\end{figure}

\section{Motivation}

A standard policy network carries one fixed set of weights, and so one fixed
answer to every opponent. When the opponent changes strategy partway through an
episode, the best a static policy can do is a compromise that is mediocre against
everyone.

If recent interaction reveals who the opponent is, the policy should be able to act
on that --- not by retraining, but by producing different weights for the opponent
in front of it.

\section{Architecture}

A context encoder reads the recent history of observations, and its output drives
the actor. The framework compares four ways of letting context shape the policy:

\begin{itemize}
  \item \resumeItem{Static baseline}{A regular, non-adaptive actor-critic
  (\texttt{static\_mlp}).}

  \item \resumeItem{Hyper head}{The context network generates the weights of the
  actor's final linear layer (\texttt{hyper\_head}).}

  \item \resumeItem{FiLM}{The context network modulates the actor's features with
  FiLM scale and shift parameters (\texttt{film}).}

  \item \resumeItem{Full hypernetwork}{The context network generates every layer of
  the actor (\texttt{full\_hyper}).}
\end{itemize}

Each is paired with a choice of context encoder: a mean over recent observations,
a flattened window, a GRU, or a small Transformer. Policies are trained with PPO.

\section{Environments}

\begin{itemize}
  \item \resumeItem{Repeated duel}{A vectorized 1v1 environment whose opponent has a
  hidden style that can switch mid-episode. Evaluation tracks reward immediately
  after a switch, during early adaptation, and once the policy has settled.}

  \item \resumeItem{Public suite}{Adapters for OpenSpiel matrix and turn-based
  games, PettingZoo rock-paper-scissors, MPE2, SlimeVolley, MAgent2 Battle, and
  Footsies.}

  \item \resumeItem{Melee Light}{Scripted opponents drawn from a pool of hidden
  styles --- rushdown, spacer, zoner, mirror, anti-frequency, and others --- and
  character matchups, switching mid-episode, so recent history identifies both
  style and matchup.}
\end{itemize}

\section{Results}

\begin{itemize}
  \item \resumeItem{Repeated duel}{In a two-seed sweep, a full hypernetwork with a
  GRU encoder reached a mean eval return of 42.3 against 33.3 for the static
  baseline. In longer same-budget runs evaluated over 256 episodes it reached a
  return of 66.9 and a 68.7\% win rate, against 58.8 and 64.7\% for the static
  baseline.}

  \item \resumeItem{Melee Light}{Against the scripted style pool, the full
  hypernetwork with a GRU encoder was the only architecture with a repeatable edge
  over the static baseline, winning 29.7\% of games against 18.8\%. Earlier runs
  against the built-in CPU showed no reliable edge for any adaptive architecture.}
\end{itemize}

\begin{figure}
\includegraphics{imgs/latent-policy-network-returns.svg}
\caption{Mean eval return in the repeated duel, two seeds, 40 updates per run.
Dark bars are adaptive architectures; the light bar is the static baseline.}
\end{figure}

\begin{figure}
\includegraphics{imgs/latent-policy-network-winrate.svg}
\caption{Win rate of the strongest adaptive architectures against the static
baseline: the repeated duel after 120 updates (256 eval episodes), and Melee Light
against the scripted style pool after 24 updates.}
\end{figure}

\end{document}
